\subsection{Одномерный поиск для BPR-функций}

Длина шага выбирается точной минимизацией
\[
    \gamma_k=\argmin_{\gamma\in[0,1]}\Psi(f^k+\gamma d^k).
\]
Для BPR-функции
\[
    \tau_e(f_e)=\bar t_e\left(1+\rho\left(\frac{f_e}{\bar f_e}\right)^{1/\mu}\right)
\]
производная одномерной функции равна
\[
    \frac{d}{d\gamma}\Psi(f^k+\gamma d^k)
    =
    \sum_{e\in\cE}
    \tau_e(f_e^k+\gamma d_e^k)d_e^k.
\]
Поэтому задача поиска $\gamma$ сводится к нахождению нуля монотонной функции на $[0,1]$.
Если производная в нуле неотрицательна, оптимум находится при $\gamma=0$.
Если производная в единице неположительна, оптимум находится при $\gamma=1$.
В остальных случаях используется одномерный численный метод.
Эта процедура одинакова для FW, CFW, BFW, FFW, WFFW и NFW; различается только направление
$d^k$.

\subsection{Итоги раздела}

Рассмотренное направление развития движется от уже известных CFW/BFW к двум новым вариантам.
NFW проверяет гипотезу, что память о большем числе прошлых направлений может усилить
эффект сопряженности.
WFFW проверяет другую гипотезу: может быть, часть ускорения объясняется не строгой
сопряженностью, а более устойчивым сглаживанием направлений.
Эксперименты показывают, что обе гипотезы полезны.
Сопряженные методы дают наиболее сильный эффект, но взвешенное сглаживание также улучшает
поведение классического FW.
